\chapter{Unterverbände und Produktkonstruktionen}
\label{b45guide:ch:constructions}

\section{Unterstrukturen und Erzeuger}

\begin{definition}[Unterverband und erzeugter Unterverband]
\label{b45guide:def:sublattice}
Eine Teilmenge $U$ eines Verbandes $L$ ist ein \emph{Unterverband},
wenn $x\wedge y,x\vee y\in U$ für alle $x,y\in U$ gilt.
Der von $X\subseteq L$ erzeugte Unterverband $\langle X\rangle_L$
ist der kleinste Unterverband, der $X$ enthält.
Leere Unterverbände sind zugelassen. Werden zusätzlich $0,1$
als Konstanten mitgeführt, muß eine Teilstruktur diese enthalten;
die entsprechende Erzeugung wird mit $\langle X\rangle_{L,0,1}$ bezeichnet.
\end{definition}

\begin{theorem}[Abschluß unter den benannten Operationen]
\label{b45guide:thm:closure}
Eine unter den Halbverbands- beziehungsweise Verbandsoperationen
abgeschlossene Teilmenge trägt mit den eingeschränkten Operationen
wieder dieselbe Struktur. Der Durchschnitt jeder Familie solcher
Unterstrukturen ist wieder eine Unterstruktur.
Für die leere Familie ist dieser Durchschnitt der ganze Träger.
Daher existieren die oben definierten erzeugten Unterstrukturen.
\cite{supclosed}
\end{theorem}

\begin{definition}[Verbandsterme]
Ein \emph{Verbandsterm} über $X$ ist ein endlicher Ausdruck, der
aus Elementen von $X$ durch die binären Zeichen $\wedge,\vee$
aufgebaut wird. Ohne benannte Konstanten besitzt jeder Term
mindestens ein Variablenvorkommen.
\end{definition}

\begin{theorem}[Terme beschreiben die gesamte Erzeugung]
\label{b45guide:thm:lattice-generation}
Für $X\subseteq L$ besteht $\langle X\rangle_L$ genau aus den
Werten aller Verbandsterme über $X$.
Mit
\[
 X_0=X,\qquad
 X_{n+1}=X_n\cup\{u\wedge v\mid u,v\in X_n\}
                 \cup\{u\vee v\mid u,v\in X_n\}
\]
gilt ausdrücklich
\[
 \langle X\rangle_L=\bigcup_{n\geq0}X_n.
\]
Jedes erzeugte Element verwendet nur endlich viele Erzeuger.
Werden $0,1$ mitgeführt, beginnt dieselbe Konstruktion bei
$X_0=X\cup\{0,1\}$. \cite{supclosed}
\end{theorem}

\begin{theorem}[Ein und zwei Erzeuger]
\label{b45guide:thm:two-generators}
In jedem Verband $L$ und für $a,b\in L$ gilt
\[
 \langle\varnothing\rangle_L=\varnothing,\qquad
 \langle\{a\}\rangle_L=\{a\},\qquad
 \langle\{a,b\}\rangle_L=\{a\wedge b,a,b,a\vee b\}.
\]
Die letzte Menge besitzt höchstens vier Elemente.
Für unvergleichbare $a,b$ besitzt sie genau vier Elemente.
Diese Aussagen beziehen sich auf die beiden binären Operationen
ohne zusätzliche benannte Konstanten.
\end{theorem}

\begin{theorem}[Verbandsintervalle]
Für $a\leq b$ in einem Verband $L$ ist
\[
 [a,b]_L=\{x\in L\mid a\leq x\leq b\}
\]
mit den eingeschränkten Operationen ein Unterverband.
Als eigenständiger beschränkter Verband hat er Null $a$ und Eins $b$.
Die Inklusion in $L$ erhält die beiden binären Operationen;
die Schranken des Intervalls müssen nicht die Schranken von $L$ sein.
\end{theorem}

\section{Produkte und Hinzufügen von Schranken}

\begin{theorem}[Koordinatenweise Produkte]
\label{b45guide:thm:products}
Sei $(S_i,\vee_i)_{i\in I}$ eine Familie von Halbverbänden.
Auf $\prod_{i\in I}S_i$ definiert
\[
 (x\vee y)_i=x_i\vee_i y_i
\]
wieder eine Halbverbandsoperation.
Für eine Familie von Verbänden erhält man durch
\[
 (x\wedge y)_i=x_i\wedge_i y_i,\qquad
 (x\vee y)_i=x_i\vee_i y_i
\]
einen Verband. Die Koordinatenprojektionen erhalten die
jeweiligen Operationen.
Das Produkt über $I=\varnothing$ ist die einelementige Struktur.
Bei benannten Schranken werden auch diese koordinatenweise gebildet.
\end{theorem}

\begin{theorem}[Äußere Null und äußere Eins]
\label{b45guide:thm:adjoin-bounds}
Jeder Halbverband $(S,\vee)$ läßt sich durch ein neues Element $0$
zu einem Halbverband mit Null erweitern, indem $0\vee x=x\vee0=x$
festgesetzt wird. Entsprechend kann durch $1\vee x=x\vee1=1$
eine neue Eins hinzugefügt werden.

Jeder Verband $L$ besitzt eine Erweiterung durch zwei neue Elemente
$0,1$, deren Operationen auf $L$ unverändert sind und
\[
 0\wedge x=0,\quad 0\vee x=x,\quad
 1\wedge x=x,\quad 1\vee x=1
\]
für alle Elemente der Erweiterung erfüllen.
Dies ergibt auch dann einen beschränkten Verband, wenn $L$ leer ist
oder bereits eigene Schranken besitzt.
\end{theorem}
